\chapter*{出版信息}
\markboth{出版信息}{出版信息}
\addcontentsline{toc}{chapter}{出版信息}

Herbert Amann

Joachim Escher

Birkhäuser Verlag Basel · Boston · Berlin

英译者：Silvio Levy 与 Matthew Cargo（Birkhäuser Verlag，2009）

作者：

Herbert Amann

Institut für Mathematik

Universität Zürich

Winterthurerstr. 190

8057 Zürich

Switzerland

电子邮箱：herbert.amann@math.uzh.ch

Joachim Escher

Institut für Angewandte Mathematik

Universität Hannover

Welfengarten 1

30167 Hannover

Germany

电子邮箱：escher@ifam.uni-hannover.de

德文原版以同一书名由瑞士 Birkhäuser Verlag 出版，© 2001 Birkhäuser Verlag

2000 数学学科分类：28-01, 28A05, 28A20, 28A25, 28B05, 58-01, 58A05, 28A10, 58C35

美国国会图书馆控制号：2008939525

Die Deutsche Bibliothek 出版书目信息

Die Deutsche Bibliothek 将本出版物列入“Deutsche Nationalbibliografie”；详细的文献数据可在互联网上通过 <http://dnb.ddb.de> 查阅。

ISBN 3-7643-7479-2 Birkhäuser Verlag, Basel – Boston – Berlin

本作品受版权保护。保留一切权利，无论是材料的全部还是部分，特别是翻译、重印、插图再利用、朗诵、广播、以缩微胶片或其他方式的复制以及存入数据库的权利。任何形式的使用均须取得版权所有人的许可。

© 2009 Birkhäuser Verlag AG，Basel · Boston · Berlin，P.O. Box 133, CH-4010 Basel, Switzerland

隶属于 Springer Science+Business Media

排版与 LaTeX：Gisela Amann，苏黎世

本书采用由无氯纸浆制成的无酸纸印刷。TCF ∞

德国印刷

ISBN 978-3-7643-7479-2

e-ISBN 978-3-7643-7480-8

987654321

www.birkhauser.ch

\chapter*{译者的话}
\markboth{译者的话}{译者的话}
\addcontentsline{toc}{chapter}{译者的话}

本书是 Herbert Amann 与 Joachim Escher 所著 \textit{Analysis III}（德文原版
\textit{Analysis III}，Birkhäuser Verlag 2001；英译者 Silvio Levy 与 Matthew Cargo，
Birkhäuser Verlag 2009）的中文译本。它与《分析学 I》《分析学 II》中文译本配套，
合起来覆盖 Amann--Escher 三卷本分析学导论的全部三卷。

\subsection*{关于这个译本}

这是一项\textbf{个人自发的翻译工程}，起因很简单：这套书写得好，而中文读者读它
要跨过一道语言的坎。第三卷收束全书：在数系、线性代数与一元分析（第一卷），
以及一元积分、多元微分与线积分（第二卷）之后，本卷把分析学推进到现代测度与积分的
框架——\textbf{第 IX 章 测度论基础}（σ-代数、外测度、可测集与勒贝格测度），
\textbf{第 X 章 积分理论}（可测函数、Bochner--Lebesgue 积分、收敛定理、$L^p$ 空间、
富比尼定理、卷积与傅里叶变换），以及\textbf{第 XI 章 流形与微分形式}与
\textbf{第 XII 章 流形上的积分}（体积测度、微分形式的积分与斯托克斯定理）。
这样一条完整的分析学链条，值得有一份尽可能贴近原文的中文本来对照阅读。

翻译以\textbf{英译本}为底本，同时参照德文原版的书名与体例。全书按章、节组织，
与原文逐节对应，公式、图表与习题全部保留；章节编号、定理编号（如「定理 5.18」）
与公式编号（如「(5.9)」）一律沿用原书，以便与英文版、德文版对照查阅。书末的
\textbf{参考文献保留英文原样、未作翻译}，便于按原文献索骥；原书的\textbf{索引
未收入本译本}——索引页码指向原书版次，与中译本重排后的页码无法一一对应。

\subsection*{关于本译本的生成方式（请先读这一段）}

必须把话说明白：\textbf{本译本由 AI（大语言模型）辅助完成}。原文的识别、初译、
排版代码，以及相当一部分核对工作，都是在人机协作下做的——机器出草稿，人定口径、
逐页抽看原书页图、逐轮返工。它\textbf{不是}由某位译者逐字推敲多年的职业译本，也
\textbf{没有经过出版社级别的编校}。

因此请读者务必注意：

\textbf{一、译文可能存在错误。}包括但不限于术语取舍、句子理解、公式与符号的抄录、
编号与交叉引用，以及个别地方与原书不一致的表述。

\textbf{二、不能保证与原著完全一致。}虽已尽量以原书页图逐页核对、把发现的讹误
逐条修掉，但「查过」不等于「查全」，遗漏在所难免。

\textbf{三、请勿把本译本当作学术引用、考试或教学的依据。}凡涉及定义、定理与证明的
准确表述，一律以 Amann--Escher 原著（英译本或德文原版）为准。

一句话：可以拿它当\textbf{对照阅读的拐杖}，但不能拿它当\textbf{标准答案}。

\subsection*{译名与记号}

数学术语以国内通行译法为准，重要的名词首次出现时保留英文原词。全书的术语取舍
与记号约定另有对照表，不在此处赘述。原文中以粗体标示的定义性词（如
\textbf{convex}、\textbf{compact}）在译文中同样以粗体标出。

\subsection*{关于校订}

原书是数字排版，正文与公式的准确字形以\textbf{原书页图}为准：文字层里大量数学
符号（积分号、范数、上下标、分式等）并不可靠，不能直接采信。原书本身的排印瑕疵
一律按原样保留、登记在案，不擅自「改正」；属于识别错误的，按页图修正。因此本译本
中个别地方的写法可能与读者的预期不同，那通常正是原书的样子。

\subsection*{内容问题反馈}

错误既然难免，就希望它被发现、被记下来。译者建了一份在线反馈表「内容问题反馈表」，
用来收集勘误、术语建议与排版问题——发现哪里不对，随手填一下即可：

\begin{center}
\url{https://github.com/JeaESenna/Amann-Analysis-zh/issues}
\end{center}


反馈时若能附上\textbf{章号、页码与原文对照}，处理起来会快得多。谢谢每一位愿意
回头指正的读者。

\subsection*{免责与致谢}

原著版权归 Herbert Amann、Joachim Escher 与 Birkhäuser Verlag 所有，英译版权
归 Silvio Levy 与 Matthew Cargo 所有。本译本仅供个人学习与研究使用，\textbf{不得
用于任何商业用途}。若本书对你有用，请务必去读原著——那是最好的感谢方式。

\vspace{2ex}
\begin{flushright}
Jea神纳\\
2026 年 10 月
\end{flushright}
\chapter*{前言}
\markboth{前言}{前言}
\addcontentsline{toc}{chapter}{前言}

第三卷为我们的分析学导论画上了句号：至此，我们为今后进一步学习这门学科所需的基础已经铺就。

与前两卷一样，本卷所包含的材料远远多于一门课程所能讲授的内容。因此，备课时，重要的是选取其中一个合适的子集；其余部分可以在讨论班中讲授，或者留给学生自学。若要快速了解本卷内容，可参阅目录与各章的导言。

本书也适合作为其他课程的背景材料，或者用于自学。我们希望，书中能让读者多次一窥更高深的分析学，从而唤起他们的好奇心，吸引他们进一步探索这一数学分支的优美与广度。

在撰写本卷的过程中，我们仰仗了朋友、同事、工作人员与学生极其宝贵的帮助。我们要特别感谢 Georg Prokert、Pavol Quittner、Olivier Steiger 与 Christoph Walker，他们以批判的眼光通读了全部书稿，从而帮助我们改正了错误，并使书稿得到了实质性的改进。我们还要感谢 Carlheinz Kneisel 与 Bea Wollenmann，他们同样阅读了书稿的大部分内容，并指出了各种不一致之处。

倘若没有我们那位“排版完美主义者”不可估量的努力，本卷不可能达到现在的样子：她对待 TeX 与其他软件的那份不知疲倦与耐心，不仅使最终成品，也使许多早先的版本达到了高度的完美。为此，我们向她致以最深的谢意。

最后，我们谨向 Thomas Hintermann 与 Birkhäuser 出版社致谢，感谢他们一如既往的灵活与友好合作。

苏黎世与汉诺威，2001 年 7 月

H. Amann、J. Escher

\subsection*{英译本前言}

我们再次向 Silvio Levy 与 Matt Cargo 致以深切的谢意，感谢他们对这部德文原著最后一部分所作的细致而准确的翻译。特别感谢 Birkhäuser Verlag 的 Thomas Hempfling 促成了本译本的问世，从而使我们的分析学教程如今能够面向更广泛的读者。

苏黎世与汉诺威，2009 年 1 月

H. Amann、J. Escher
